\documentclass[11pt]{report}

\usepackage[margin=1in]{geometry}
\usepackage{hyperref}
\usepackage{listings}
\usepackage{xcolor}
\usepackage{parskip}
\usepackage{enumitem}
\usepackage{longtable}
\usepackage{booktabs}
\usepackage{seqsplit}
\usepackage{tocloft}

\lstdefinelanguage{bmake}{
  morecomment=[l]{\#},
  morestring=[b]",
  sensitive=true
}
\lstset{
  basicstyle=\ttfamily\small,
  breaklines=true,
  frame=single,
  framesep=3pt,
  columns=fullflexible,
  backgroundcolor=\color{black!3}
}

\hypersetup{
  colorlinks=true,
  linkcolor=blue!50!black,
  urlcolor=blue!50!black
}

\title{\textbf{Bmake It}\\[0.3em]\large A BSD-Make-Based Cross-Platform Build System\\[0.3em]\normalsize Package Documentation}
\author{Draft v1}
\date{\today}

\newcommand{\reqfoot}[1]{\footnote{\footnotesize\seqsplit{#1}}}

\begin{document}

\maketitle

\begin{abstract}
\noindent
This document describes \texttt{Bmake It}, a BSD-make-based build
system for C, C++, and lex/yacc grammar sources, designed to mimic the
conceptual model of Dassault Systèmes' RADE/\texttt{mkmk} tool
(workspaces, frameworks, modules, prerequisites) while reusing real BSD
make infrastructure wherever it fits\reqfoot{CON-use-bsdmake}. Every
substantive claim in this document is footnoted with the requirement,
decision, or constraint alias from the project's discovery-driven
design record (\texttt{project-model/}) that established it, so any
statement here can be traced back to its origin and rationale.

This is a \emph{draft} document, written while the design is still
being finalized; sections marked \textbf{[open]} describe areas where
the underlying design is intentionally still undecided or deferred.
\end{abstract}

\tableofcontents

%===============================================================
\chapter{Overview}
%===============================================================

\section{What Bmake It Is}

\texttt{Bmake It} is a cross-platform build system for C and C++
(plus lex/yacc grammar sources, which compile down to C/C++)\reqfoot{REQ-supported-languages},
built on top of real BSD make rather than GNU make or a bespoke build
tool. Its object model — \emph{workspaces} containing \emph{frameworks}
containing \emph{modules} — is deliberately modeled on Dassault
Systèmes' RADE/\texttt{mkmk} tool, a CAA (Component Application
Architecture) build system used internally at Dassault Systèmes since
the late 1990s. Where \texttt{mkmk}'s specific mechanisms don't fit
BSD make's own idioms well, \texttt{Bmake It} diverges deliberately
rather than reproducing them exactly\reqfoot{CON-use-bsdmake}.

Three mechanisms present in real \texttt{mkmk} are explicitly
\emph{not} carried over, by design rather than oversight:

\begin{itemize}
\item No dedicated ``identity card'' special file for declaring
  prerequisites — prerequisites are declared directly in the ordinary
  framework \verb|makefile|\reqfoot{NREQ-no-identitycard-special-file}.
\item No four-level RADE header hierarchy (Public / Protected /
  Private / Local) — \texttt{Bmake It} uses a simpler three-level
  scheme instead, described in Chapter~\ref{ch:basic-usage}\reqfoot{NREQ-no-four-level-header-visibility}.
\item No \texttt{mkGetPreq}/\texttt{mkCopyPreq}-style commands for
  managing prerequisite search paths — prerequisite resolution is
  purely dynamic, with no static-copy mode and no dedicated setup
  command\reqfoot{NREQ-no-copy-preq-commands}.
\end{itemize}

\section{Design Philosophy}

The guiding principle throughout is: \emph{reuse real BSD make
variables and mechanisms wherever they already exist; introduce new,
project-specific ones only where BSD make genuinely has no
equivalent.} This shows up consistently:

\begin{itemize}
\item Target selection reuses BSD make's own \verb|TARGET=|/\verb|TARGET_ARCH=|
  convention rather than inventing new names\reqfoot{REQ-target-arch-toolchain-selection-is-custom-req}.
\item Installation reuses \verb|DESTDIR| unchanged; \verb|PREFIX|, it
  turns out, is not actually a native BSD make concept at all (see
  Section~\ref{sec:validation})\reqfoot{REQ-prefix-is-new-not-reused-req}.
\item Compile and link flags (\verb|CFLAGS|, \verb|CXXFLAGS|,
  \verb|LDFLAGS|, \verb|YFLAGS|, \verb|LFLAGS|) are the ordinary BSD
  make variables, extended via \verb|+=| — no wrapper macros the way
  real \texttt{mkmk} used \verb|LOCAL_CFLAGS|-style
  names\reqfoot{REQ-flags-reuse-bsd-native-vars-req}.
\item Debug-vs-release is controlled via BSD make's own
  \verb|DEBUG_FLAGS|, not a new \verb!MODE=debug|release! abstraction\reqfoot{REQ-debug-release-via-native-bsd-flags-req}.
\item Shared-library versioning reuses \verb|SHLIB_MAJOR|/\verb|SHLIB_MINOR|
  directly from real \verb|bsd.lib.mk|\reqfoot{REQ-shlib-major-minor-cross-platform-emission-req}.
\end{itemize}

New, project-specific concepts (\verb|PARENT_WS|, \verb|PREREQS|,
\verb|PROG|, \verb|LIB|, \verb|LIBS|, \verb|LIB_SHARED|, \verb|INCL|,
\verb|TOOLCHAIN|) exist only for domain concepts BSD make has no
native equivalent for at all — workspace parenting, framework
prerequisites, module identity, linking, header promotion, and
toolchain bundle selection\reqfoot{REQ-uniform-makefile-filename-req}.

\textbf{This reuse-first philosophy has real, tested limits.} An
empirical validation pass against real FreeBSD \verb|share/mk| and a
portable \texttt{bmake} binary found that several assumed
``native reuse'' cases were actually wrong: \verb|LIB_SHARED| has no
native BSD-make equivalent at all (real \verb|bsd.lib.mk| builds a
static archive whenever \verb|LIB| is set, and \emph{additionally}
builds a shared library only if \verb|SHLIB_MAJOR| is
present)\reqfoot{REQ-lib-shared-is-new-not-reused-req}, \verb|PREFIX|
does not exist in \verb|bsd.prog.mk|/\verb|bsd.lib.mk| at
all\reqfoot{REQ-prefix-is-new-not-reused-req}, and
\verb|TARGET|/\verb|TARGET_ARCH| have no effect whatsoever without
FreeBSD's full base-system build machinery, which a standalone project
does not have\reqfoot{REQ-target-arch-toolchain-selection-is-custom-req}.
These corrections are described in Section~\ref{sec:validation} and
are already reflected in the macro reference in
Chapter~\ref{ch:advanced}.

\section{The Domain Model}
\label{sec:domain-model}

\begin{description}
\item[Workspace] A set of frameworks\reqfoot{REQ-workspace-is-set-of-frameworks}.
  A workspace may declare an ordered list of \emph{parent} workspaces;
  the build searches them, in declared order, for anything not found
  locally\reqfoot{REQ-workspace-parent-search-order}. This search is
  purely dynamic — there is no physical-copy mode and no separate
  setup command, unlike real \texttt{mkmk}'s
  \texttt{mkGetPreq}/\texttt{mkCopyPreq}\reqfoot{DEC-prereq-access-dynamic-only}.
  A workspace carries no resource files of its
  own\reqfoot{NREQ-no-workspace-level-share-source}.

\item[Framework] Contains one or more modules, exposes public headers,
  and declares its prerequisite frameworks\reqfoot{REQ-framework-contains-modules}.
  A framework's prerequisite declaration is \emph{mandatory}, even when
  empty — its presence is what makes a directory auto-discoverable as
  a framework in the first place\reqfoot{REQ-prereqs-mandatory-even-empty-req}.

\item[Module] The basic link unit: a module produces exactly one
  executable, static library, or shared
  library\reqfoot{REQ-module-is-basic-link-unit}. A module's directory
  is named \verb|NAME.m|\reqfoot{REQ-module-dir-naming-convention}, and
  modules are auto-discovered by scanning — no explicit list is
  maintained by hand\reqfoot{REQ-modules-and-frameworks-auto-discovered-req}.

\item[Resource files] Non-source runtime data (icons, configuration
  files, and similar), carried by both modules and
  frameworks\reqfoot{REQ-resource-file-model}, staged through a
  three-layer \verb|share/| overlay described in
  Section~\ref{sec:resources}.
\end{description}

\section{Platform and Language Scope}

\textbf{Languages}: C, C++, and lex/yacc grammar sources. Java,
Fortran, IDL, and Express — all supported by real \texttt{mkmk} — are
explicitly out of scope for now\reqfoot{REQ-supported-languages}.

\textbf{Hosts}: macOS (via MacPorts), Linux, and the \texttt{*BSD}
family\reqfoot{CON-cross-platform-targets}, plus Cygwin and WSL treated
as \texttt{*NIX}-like build hosts (WSL is, for build purposes, simply
Linux; no native MSVC toolchain is required)\reqfoot{REQ-windows-hosts-cygwin-wsl}.

\textbf{Phase 1} — the concrete, currently-targeted scope — is exactly
four build targets\reqfoot{REQ-phase1-concrete-targets}:
\begin{center}
\begin{tabular}{ll}
\toprule
Target key & Meaning \\
\midrule
\verb|macos-arm64| & macOS, Apple Silicon, LLVM (default toolchain) \\
\verb|macos-arm64-gcc| & macOS, Apple Silicon, GCC \\
\verb|freebsd-amd64| & FreeBSD, x86-64, LLVM (default toolchain) \\
\verb|freebsd-amd64-gcc| & FreeBSD, x86-64, GCC \\
\bottomrule
\end{tabular}
\end{center}

Everything outside this set — OpenBSD/NetBSD, per-distribution Linux
labeling, kernel-level or Darwin/XNU development, macOS \texttt{.app}
bundle output, the exact \texttt{win} ABI story (MinGW vs.\ Cygwin vs.\
MSVC), musl libc — is deliberately deferred, not designed against yet\reqfoot{REQ-userland-only-scope}.

%===============================================================
\chapter{Basic Usage}
\label{ch:basic-usage}
%===============================================================

This chapter covers the minimum needed to create a workspace, populate
it with a framework and a module, and build it.

\section{Directory Layout at a Glance}

Every level of the hierarchy — workspace, framework, module — uses a
file literally named \verb|makefile| (not \verb|makefile.mk|), which
\verb|.include|s a role-specific helper file, mirroring the way real
BSD make's own \verb|bsd.prog.mk|/\verb|bsd.lib.mk| work\reqfoot{REQ-uniform-makefile-filename-req}:

\begin{center}
\begin{tabular}{lll}
\toprule
Level & Makefile contents (first line) & Key macro(s) \\
\midrule
Workspace & \verb|.include <mk.workspace.mk>| & \verb|PARENT_WS=| \\
Framework & \verb|.include <mk.framework.mk>| & \verb|PREREQS=| \\
Module (executable) & \verb|.include <mk.prog.mk>| & \verb|PROG=|, \verb|LIBS=| \\
Module (library) & \verb|.include <mk.lib.mk>| & \verb|LIB=|, \verb|LIB_SHARED=|, \verb|LIBS=| \\
\bottomrule
\end{tabular}
\end{center}

\section{Creating a Workspace}

A workspace is just a directory containing a \verb|makefile|:

\begin{lstlisting}[language=bmake]
.include <mk.workspace.mk>
PARENT_WS=/abs/path/to/parent-ws-1 /abs/path/to/parent-ws-2
\end{lstlisting}

\verb|PARENT_WS| is an ordered, space-delimited list of \textbf{absolute}
paths to parent workspaces\reqfoot{REQ-workspace-parent-macro-absolute-req}
— relative paths were explicitly considered and rejected as too
fragile. If left empty, the workspace has no parents and everything it
needs must live inside it. Frameworks inside the workspace are
auto-discovered: any immediate subdirectory containing a \verb|makefile|
that declares \verb|PREREQS=| is treated as a
framework\reqfoot{REQ-framework-discovery-signal-req}; there is no
explicit list to maintain.

\section{Creating a Framework}

A framework is a directory (any name — no special suffix, unlike
modules) containing its own \verb|makefile|:

\begin{lstlisting}[language=bmake]
.include <mk.framework.mk>
PREREQS=System
\end{lstlisting}

\verb|PREREQS=| is \textbf{mandatory even when empty}
(\verb|PREREQS=| with nothing after it is valid and means ``no
prerequisites'')\reqfoot{REQ-prereqs-mandatory-even-empty-req} — this
is not merely documentation, it is the literal signal the workspace
uses to recognize the directory as a framework at all. The list is
ordered, all entries are implicitly public (there is no per-entry
visibility qualifier the way real \texttt{mkmk}'s identity card
distinguished Public/Protected/Private), and \verb|Bmake It| does
\emph{not} validate the ordering or detect
cycles\reqfoot{REQ-framework-prereqs-macro-req}: an incorrectly
ordered or cyclic \verb|PREREQS| list simply surfaces later as an
ordinary ``header not found'' compile
error\reqfoot{REQ-prereqs-misorder-no-validation-req}.

A framework has three directories for hand-written headers,
described fully in Section~\ref{sec:headers}, and an optional
\verb|share/| directory for resources (Section~\ref{sec:resources}).

\section{Creating a Module}

A module's directory is always named \verb|NAME.m|; \verb|NAME| also
becomes the default output name if not overridden\reqfoot{REQ-module-dir-naming-convention, REQ-prog-lib-default-to-module-name-req}.
Whether \verb|NAME.m| and \verb|name.m| are distinct depends entirely
on the host filesystem's own case sensitivity — \texttt{Bmake It}
does not itself enforce or normalize case\reqfoot{REQ-module-name-case-sensitivity-req}.

\subsection*{An executable module}
\begin{lstlisting}[language=bmake]
.include <mk.prog.mk>
PROG=hello        # optional -- defaults to the directory name
LIBS=greet        # link-time only, space-delimited
\end{lstlisting}

\subsection*{A library module}
\begin{lstlisting}[language=bmake]
.include <mk.lib.mk>
LIB=greet             # optional -- defaults to the directory name
LIB_SHARED=YES        # default; LIB_SHARED=NO builds a static archive instead
LIBS=
\end{lstlisting}

\verb|LIB_SHARED| defaults to \verb|YES|;
setting it explicitly to \verb|NO| builds a static/archive library
instead\reqfoot{REQ-lib-shared-is-new-not-reused-req}. \textbf{Important, and
confirmed by empirical testing against real \texttt{bsd.lib.mk}}: this
\verb|YES|/\verb|NO| toggle is not itself a native BSD make mechanism
— real \verb|bsd.lib.mk| unconditionally builds a static archive
whenever \verb|LIB| is set, and additionally builds a shared library
only when \verb|SHLIB_MAJOR| is also present. \texttt{Bmake It}'s
\verb|mk.lib.mk| implements the \verb|LIB_SHARED| abstraction on top of
that native behavior\reqfoot{REQ-lib-shared-is-new-not-reused-req}.

A module's \verb|src/| directory is scanned automatically: every
source file physically present there is treated as compilation input,
with no \verb|SRCS=|-style list to maintain by
hand\reqfoot{REQ-module-src-auto-discovery}. \verb|.y| and \verb|.l|
(yacc and lex) files are compiled through the same automatic
mechanism, confirmed working via built-in suffix transformation rules
with zero explicit rule declarations\reqfoot{OBS-validation-yacc-lex-suffix-rules-confirmed}.

\verb|LIBS=| lists other modules or system libraries this module links
against. It is \textbf{purely a link-time instruction} and has no
effect whatsoever on header visibility, even for sibling modules in
the same framework\reqfoot{REQ-libs-macro-link-only-headers-via-prereqs-req}
— see Section~\ref{sec:headers} for how header visibility actually
works.

\section{Header Visibility}
\label{sec:headers}

Headers have exactly three visibility levels\reqfoot{REQ-three-level-header-visibility-req}:

\begin{center}
\begin{tabular}{lll}
\toprule
Level & Directory & Visible to \\
\midrule
Public & \verb|fw/include/| & Any framework declaring this one in \verb|PREREQS=| \\
Framework-internal & \verb|fw/local/include/| & All of this framework's own modules \\
Module-private & \verb|module.m/include/| & Only that module's own source files \\
\bottomrule
\end{tabular}
\end{center}

Header visibility is resolved \textbf{exclusively} through the
containing framework's \verb|PREREQS=| declaration — never through a
module's \verb|LIBS=|\reqfoot{REQ-headers-resolved-from-source-req}.
This is a genuine, deliberate design point worth internalizing: a
module can link against a library via \verb|LIBS=| while remaining
completely unable to see that library's headers, if the containing
frameworks aren't in a \verb|PREREQS=| relationship. Headers are
always read directly from source (\verb|include/| directories) —
never from build output — because headers are plain text and don't
need per-target staging the way compiled artifacts
do\reqfoot{REQ-headers-resolved-from-source-req}. The one exception is
\emph{generated} headers (from yacc/lex), covered in
Section~\ref{sec:generated-headers}.

\section{Building}

Run \verb|make| in a module directory to build just that module, in a
framework directory to build all its modules, or at the workspace root
to build everything. Modules and frameworks are auto-discovered by
scanning — modules via the \verb|.m| suffix, frameworks via the
\verb|makefile|+\verb|PREREQS=| signal\reqfoot{REQ-modules-and-frameworks-auto-discovered-req}.

The build strategy is \emph{fused}: each module compiles, generates
headers if applicable, and links in one pass — matching real
\verb|bsd.*.mk| idiom directly, rather than a global
compile-everything-then-link-everything
split\reqfoot{REQ-fused-build-per-module-full-dependency-order-req}.
Build order across frameworks follows the \verb|PREREQS=| dependency
graph; build order across sibling modules follows the \verb|LIBS=|
graph plus any consumed generated-header
dependency\reqfoot{REQ-framework-build-order-follows-prereqs-graph-req, REQ-fused-build-per-module-full-dependency-order-req}.
The concrete mechanism generating this order — confirmed by a working
prototype, not merely designed on paper — is described in
Section~\ref{sec:dep-order}.

\begin{lstlisting}[language=bash]
$ cd myworkspace/Hello
$ make                    # builds every module in Hello, in dependency order
$ cd hello.m
$ make                    # builds just this one module
\end{lstlisting}

\section{Cleaning and Installing}

\verb|make clean| removes \textbf{every} generated artifact for the
selected target — not just the compiled binaries, but also the
\verb|distrib/| packaging staging area, per-module \verb|.depend|
files, and the internal dependency-resolution
caches\reqfoot{REQ-clean-removes-all-derived-artifacts-req}. After
\verb|make clean TARGET=all|, the tree should be indistinguishable
from a fresh checkout.

\verb|make install DESTDIR=<staging-root> PREFIX=<final-path>| copies
the built \verb|bin|/\verb|lib|/\verb|share| content into
\verb|$(DESTDIR)$(PREFIX)|. \verb|DESTDIR| is confirmed genuinely
native to BSD make. \verb|PREFIX|, however, is \textbf{not} — testing
found zero references to it anywhere in real
\verb|bsd.own.mk|/\verb|bsd.prog.mk|/\verb|bsd.lib.mk|; real FreeBSD
base-system installs use hardcoded individual \verb|*DIR| variables
(\verb|BINDIR|, \verb|LIBDIR|, etc.) directly rather than a single
unifying \verb|PREFIX| the way GNU autotools or MacPorts do it.
\texttt{Bmake It}'s \verb|mk.prog.mk|/\verb|mk.lib.mk| derive those
\verb|*DIR| variables from \verb|PREFIX=| themselves\reqfoot{REQ-prefix-is-new-not-reused-req}.
It requires the target to already be built — it will not implicitly
trigger a build for you\reqfoot{REQ-install-requires-prebuilt-target-req}.
\textbf{Note}: confirmed by testing against both real FreeBSD and
NetBSD's own \verb|share/mk|, BSD make's \verb|install| does not
auto-create missing \verb|DESTDIR| parent
directories\reqfoot{OBS-validation-destdir-native-prefix-is-not} — the
destination tree must already exist.

%===============================================================
\chapter{Advanced Features}
\label{ch:advanced}
%===============================================================

\section{Target Selection and Cross-Compilation}

The full target key has the shape
\verb|<os>-<arch>[-<toolchain>][-<abi>]|, with each bracketed segment
omitted whenever it equals that platform's default\reqfoot{REQ-target-key-omit-defaults-req}:

\begin{center}
\begin{tabular}{lll}
\toprule
Segment & Default (omitted) & Shown when non-default \\
\midrule
Toolchain & \verb|llvm| & \verb|gcc| \\
ABI/libc & \verb|msvc| (win); \verb|glibc| (linux) & \verb|mingw|, \verb|cygwin|; \verb|musl| \\
\bottomrule
\end{tabular}
\end{center}

On the command line:
\begin{lstlisting}[language=bash]
$ make TARGET=freebsd TARGET_ARCH=amd64 TOOLCHAIN=gcc
\end{lstlisting}

\verb|TARGET=| and \verb|TARGET_ARCH=| reuse real BSD make's own
FreeBSD cross-build variable names\reqfoot{REQ-target-arch-toolchain-selection-is-custom-req};
\verb|TOOLCHAIN=| is new, since BSD make has no native concept of a
toolchain \emph{bundle} (compiler, linker, and archiver selected
together)\reqfoot{REQ-toolchain-cli-variable-req}. Setting
\verb|TOOLCHAIN=| \verb|.include|s a dedicated
\verb|mk.toolchain.<name>.mk|, which sets \verb|CC|/\verb|CXX|/\verb|LD|/\verb|AR|/\verb|AS|
as a coherent bundle; these remain individually overridable afterward
for mixed setups (e.g.\ Clang with GNU \verb|ld|)\reqfoot{REQ-toolchain-triggers-dedicated-mk-file-req}.

\textbf{A significant, empirically-confirmed correction}: unlike most
of the variables reused directly from BSD make,
\verb|TARGET|/\verb|TARGET_ARCH| do \emph{not} come with working
cross-compilation ``for free.'' Testing against real FreeBSD
\verb|share/mk| showed that setting \verb|TARGET_ARCH=arm64| on an
\texttt{amd64} host still produced a plain \texttt{x86-64} binary —
these variables are inert without FreeBSD's base-system
\verb|Makefile.inc1|/\verb|buildworld| machinery, which a standalone
project does not have\reqfoot{REQ-target-arch-toolchain-selection-is-custom-req}.
Further investigation across the whole BSD-make family (FreeBSD,
NetBSD, and generic portable \texttt{pmake}) found the identical
pattern everywhere: real cross-compilation always happens via an
\emph{external orchestration layer}, with \verb|TARGET|/\verb|TARGET_ARCH|-style
variables used only as \emph{data} for architecture-conditional
branching, never as a trigger that itself selects a compiler binary.
Consequently, \verb|mk.toolchain.<name>.mk| must itself map
\verb|(TARGET, TARGET_ARCH, TOOLCHAIN)| combinations to concrete
compiler paths\reqfoot{REQ-toolchain-mk-must-set-explicit-cc-paths-req, OBS-validation-cross-compile-is-external-orchestration-everywhere}.

\section{Resource Files}
\label{sec:resources}

Both modules and frameworks may carry a \verb|share/| directory for
non-source runtime data\reqfoot{REQ-resource-file-model}. Resources
support per-target variants via a three-level fallback hierarchy:

\begin{center}
\begin{tabular}{ll}
\toprule
Directory & Applies to \\
\midrule
\verb|share/common/| & Every target, unconditionally \\
\verb|share/<os>/| & All architectures of one OS (e.g.\ \verb|share/freebsd/|) \\
\verb|share/<os>_<arch>/| & One specific OS+architecture pairing \\
\bottomrule
\end{tabular}
\end{center}

Staging is a \textbf{layered overlay copy}, not a
pick-one-directory selection: \verb|common/| copies first, then
\verb|<os>/| on top of it (overwriting same-named files), then
\verb|<os>_<arch>/| on top of that\reqfoot{REQ-share-overlay-copy-order}.
Files unique to any one layer survive regardless; only files present
at multiple layers end up as the most specific version. Unlike the
compiled-artifact target key, the resource fallback key is
\textbf{OS+architecture only} — toolchain never affects which resource
variant is selected\reqfoot{REQ-share-dir-fallback-hierarchy}.

\section{Generated Headers}
\label{sec:generated-headers}

Headers produced by tooling (yacc/lex output, for
example) don't exist until the build runs, so they can't live in
\verb|include/| as source. They are, by default, \textbf{private} to
the generating module, landing in
\verb|module/build/<KEY>/include/|\reqfoot{REQ-generated-headers-in-build-tree-req}.

A module promotes specific generated headers to public visibility via
\verb|INCL=|:
\begin{lstlisting}[language=bmake]
INCL=grammar.h
\end{lstlisting}
On build, the listed headers are copied up into
\verb|fw/build/<KEY>/include/| — a framework-level directory
containing \textbf{only} promoted generated headers, never hand-written
ones\reqfoot{REQ-incl-macro-promotes-generated-headers-req}. When
another framework resolves a \verb|PREREQS|-declared framework's
headers, it checks \emph{both} that framework's \verb|include/|
(hand-written) \emph{and} its \verb|build/<KEY>/include/| (promoted
generated) — the one deliberate, scoped exception to ``never consult
build output for headers''\reqfoot{REQ-generated-headers-in-build-tree-req}.
There is no workspace-level equivalent: a workspace-wide generated-header
directory is never consulted at compile time, only during packaging.

\section{Versioned Shared Libraries}

Shared-library versioning reuses \verb|SHLIB_MAJOR|/\verb|SHLIB_MINOR|
directly, translated per platform by \verb|mk.lib.mk|:

\begin{center}
\begin{tabular}{ll}
\toprule
Platform & Emission \\
\midrule
FreeBSD / Linux (ELF) & \verb|libfoo.so.MAJOR[.MINOR]| + standard symlink chain \\
macOS (Mach-O) & \verb|libfoo.MAJOR.dylib| + symlink + linker \verb|-compatibility_version|/\verb|-current_version| \\
\bottomrule
\end{tabular}
\end{center}

A single author-facing macro pair drives both platform-appropriate
emission strategies, since embedded Mach-O version metadata (not just
the filename) is macOS's actual ABI-compatibility
mechanism\reqfoot{REQ-shlib-major-minor-cross-platform-emission-req}.
This specific translation for macOS remains \textbf{unverified against
real MacPorts \texttt{bsdmake}} — no macOS environment was available
during this project's validation pass, and real \verb|bsd.lib.mk|'s
shared-library logic was designed primarily for ELF; MacPorts'
\verb|bsdmake| may have no Mach-O awareness at all, in which case this
logic needs to be written entirely inside \texttt{Bmake It}'s own
\verb|mk.lib.mk|. \textbf{[open]}

\section{Packaging}

Packaging is driven by separate named \verb|make| targets, one per
format, rather than a single parameterized
target\reqfoot{REQ-packaging-per-format-targets-req}: confirmed target
names so far are \verb|port| (BSD ports-style source recipe) and
\verb|pkg| (macOS-style binary package), with \verb|deb|/\verb|rpm|
following the same pattern. Every binary-package target (\verb|pkg|,
\verb|deb|, \verb|rpm|) stages a \verb|distrib/<KEY>/| directory — the
same target-key scheme as \verb|build/| — by internally invoking the
equivalent of \verb|install DESTDIR=distrib/<KEY>|, reusing the install
machinery rather than a separate copy
mechanism\reqfoot{REQ-packaging-invokes-install-into-distrib-req}.

\verb|make port| is \textbf{categorically different}: it produces a
source-based, BSD-ports-style recipe (port \verb|Makefile|,
\verb|distinfo|, patches, \verb|pkg-plist|) for the ports system to
fetch and build itself — no pre-built binary is ever staged or
embedded\reqfoot{REQ-port-target-is-source-recipe-not-binary-package-req}.
Its exact generation mechanism is deferred to a later
phase\reqfoot{REQ-make-port-deferred-req}. \textbf{[open]}

\section{Running a Built Program}

\verb|make run PROGRAM=<name> ARGS='<args>'|, run at workspace level,
sets the platform-appropriate dynamic-library search-path variable
(\verb|LD_LIBRARY_PATH| on Linux/FreeBSD, \verb|DYLD_LIBRARY_PATH| on
macOS) to include the workspace's consolidated \verb|build/<KEY>/lib/|,
then executes \verb|build/<KEY>/bin/<PROGRAM>| with the given
arguments\reqfoot{REQ-run-target-req} — mirroring real
\texttt{mkmk}'s own \texttt{mkrun} command, adapted to
\texttt{Bmake It}'s own \verb|build/<KEY>/| layout.

\section{Self-Documentation: \texttt{make help}}

\verb|make help|, run at any level, lists every valid target and
variable applicable there, along with a description and — where
relevant — the set of valid values (for example,
\verb!LIB_SHARED=YES|NO!). Each role-specific \verb|mk.*.mk| file is
responsible for documenting its own contribution, so the aggregate
output always matches exactly what's actually
\verb|.include|d at that level\reqfoot{REQ-help-target-req}.

\section{Convenience Targets}

Two small convenience targets exist purely to reduce hand-editing
error, without introducing any resolution or search behavior (they do
\textbf{not} reopen the earlier rejection of
\texttt{mkGetPreq}/\texttt{mkCopyPreq}):

\begin{lstlisting}[language=bash]
$ make add-prereq FW=System     # appends System to PREREQS= in the current framework
$ make add-parent WS=/abs/path  # appends /abs/path to PARENT_WS=
\end{lstlisting}
\reqfoot{REQ-add-prereq-add-parent-convenience-targets-req}

\section{Dependency-Order Generation}
\label{sec:dep-order}

Since real \verb|bsd.subdir.mk| has no concept of a topological build
order — it simply processes whatever \verb|SUBDIR| list it's given —
\texttt{Bmake It} generates that order itself, in a mechanism
confirmed by a working, tested prototype rather than only designed on
paper\reqfoot{REQ-dep-order-generation-mechanism-req}:

\begin{enumerate}
\item Discover \verb|.m| module directories by scanning.
\item Query each module's \verb|LIBS=| \emph{in isolation}, via
  \verb|bmake -f <dir>/makefile -V LIBS| — confirmed to read the
  variable without triggering a build.
\item Topologically sort (Kahn's algorithm), silently ignoring any
  \verb|LIBS=| entry that isn't a locally-discovered module (such
  references are resolved via the framework's \verb|PREREQS=|
  instead, not module-level ordering).
\item Emit the sorted order to a generated \verb|.mk| fragment
  (\verb|SUBDIR=modA.m modB.m ...|), \verb|.include|d before
  \verb|.include <bsd.subdir.mk>|.
\end{enumerate}

This was tested on a real four-module dependency chain
(\texttt{A}$\leftarrow$\texttt{B}$\leftarrow$\texttt{C}$\leftarrow$\texttt{D},
with \texttt{C} depending on both \texttt{A} and \texttt{B}) and built
in the correct order every time; a genuine cycle
(\texttt{X}$\leftrightarrow$\texttt{Y}) was confirmed to fail loudly
with a clear diagnostic rather than hanging or silently producing a
broken order. The identical mechanism generalizes directly to
workspace-level framework ordering, querying \verb|PREREQS=| instead
of \verb|LIBS=|.

\subsection*{The \texttt{order} Target}
\label{sec:order-target}

This generation step is formalized as its own explicit,
workspace-level target: \verb|make order| performs a \textbf{blank
run} — no compilation, no linking — scanning every framework and every
module across the whole workspace, resolving the complete dependency
graph (framework order via \verb|PREREQS=|, and each framework's
module order via \verb|LIBS=|) in one pass, and writing the result to
a cache\reqfoot{REQ-order-target-blank-run-req}. The default
\verb|make| target then reads that cache rather than re-resolving
order inline on every invocation.

The order cache is treated as an \emph{implicit prerequisite} of the
default build, not an optional manual step: if it is missing or
detected stale, \verb|make| regenerates it automatically before
proceeding, so a fresh checkout builds correctly with plain
\verb|make| and no separate mandatory
invocation\reqfoot{REQ-order-cache-auto-triggered-req}. \verb|make order|
remains independently runnable on its own, for inspecting the
resolved build order without triggering an actual build.

As a defensive safety net on top of this — since the order cache
\emph{should} always be correct, but a stale cache, a partial manual
build, or a genuine bug could in principle defeat it — immediately
before invoking the linker for a module, \texttt{Bmake It} verifies
that every \verb|LIBS=| entry's build artifact actually exists. A
missing one produces an explicit, specific error (\emph{``cannot link
\texttt{hello}: prerequisite library \texttt{greet} (from
\texttt{LIBS=}) has not been built yet''}) rather than the linker's
own generic \emph{``cannot find -lgreet''}\reqfoot{REQ-link-time-not-yet-built-guard-req}.
This should essentially never fire in ordinary operation; its value is
a clear diagnostic on the rare occasion something has genuinely gone
wrong.

\section{Caching}

Two distinct caches, both modeled on real BSD make's own
\verb|.depend| idiom (a generated file, regenerated via an explicit or
detected-stale trigger, consumed via \verb|.include| on subsequent
runs)\reqfoot{REQ-prereqs-resolution-cached-like-depend-req}:

\begin{enumerate}
\item \textbf{\texttt{PARENT\_WS}/\texttt{PREREQS} resolution cache} —
  workspace/framework-level: caches header search paths and build
  order so they aren't re-walked on every invocation. Formalized as an
  explicit workspace-level target, \verb|make order| — described in
  Section~\ref{sec:order-target}.
\item \textbf{Per-module \texttt{.depend}} — file-level,
  \texttt{mkdep}-style: which source files include which headers
  (transitively), so touching a header only triggers recompilation of
  the affected \verb|.o| files, not a full module rebuild.
\end{enumerate}

\textbf{The per-module case has a genuinely interesting empirical
story, worth telling in full rather than just stating the conclusion.}
An initial test ran the separate \verb|make depend| target against
real FreeBSD \verb|share/mk| and found it nearly useless — no header
tracking, just one link-level line. That looked like confirmation the
mechanism needed custom shell-script wiring. A closer look reversed
that conclusion: every ordinary compile already includes
\verb|-MD -MF.depend.<obj>| — \verb|bsd.sys.mk| wires the
\emph{compiler's own} native dependency-generation flag into every
build, with no separate step at all, and this was confirmed correct
end-to-end (touching a shared header triggered recompilation of
exactly the right objects). But a further cross-check against
NetBSD's own bundled \verb|mk| set found \emph{zero} such wiring there
— plain compiles, no dependency flags, no rebuild when a header
changed. So the working mechanism is real, but specific to FreeBSD's
own actual \verb|bsd.sys.mk|, not a general BSD-make-family property.
Since macOS (via MacPorts) is a phase-1 host and MacPorts' own
\verb|bsdmake| could not be tested in this pass, whether it includes
equivalent wiring is genuinely unknown\reqfoot{REQ-md-dependency-tracking-freebsd-specific-req}
— though the fallback is simple either way: \verb|mk.prog.mk|/\verb|mk.lib.mk|
can add the same \verb|-MD| flags themselves if MacPorts turns out to
lack them.

Exact generated filenames, staleness-detection mechanisms, and
regeneration triggers for the \verb|PARENT_WS|/\verb|PREREQS| cache
are deliberately deferred until real \verb|bsd.*.mk| infrastructure
can be checked hands-on for what it already
provides\reqfoot{OBS-cache-mechanism-needs-empirical-check}. \textbf{[open]}

\section{Empirical Validation Summary}
\label{sec:validation}

A number of claims in this document rest on empirical testing rather
than documentation alone — real \verb|bsd.*.mk| behavior sometimes
diverges from reasonable-sounding assumptions, and more than once in
this project's own validation pass, a first test result turned out to
need correcting after a closer look. Testing was performed using the
portable \texttt{bmake} package (a NetBSD-derived,
FreeBSD/MacPorts-lineage-compatible \texttt{make}) against both a real
FreeBSD \verb|share/mk| checkout and NetBSD's own bundled \verb|mk|
set\reqfoot{OBS-validation-yacc-lex-suffix-rules-confirmed}. A full
account, organized by risk tier — including the corrections themselves,
not just their final conclusions — is maintained separately as
\texttt{spec/60-bsd-make-validation-plan.md}; the headline results:

\begin{center}
\begin{longtable}{p{0.3\textwidth}p{0.6\textwidth}}
\toprule
Assumption tested & Result \\
\midrule
\endhead
\verb|.y|/\verb|.l| suffix rules & Confirmed working automatically, zero explicit rules needed \\
\verb|DESTDIR| staging & Confirmed genuinely native, works as expected \\
\verb|DEBUG_FLAGS| & Confirmed genuinely native, real working mechanism \\
Dynamic \verb|SUBDIR!=| population & Confirmed working — the whole auto-discovery design rests on this \\
\verb!LIB_SHARED=YES|NO! & \textbf{Not native} — real \verb|bsd.lib.mk| has no such toggle \\
\verb|TARGET|/\verb|TARGET_ARCH| & \textbf{Inert} without FreeBSD base-system machinery \\
\verb|PREFIX| & \textbf{Not native} — does not exist in \verb|bsd.prog.mk|/\verb|bsd.lib.mk| \\
Per-object \verb|-MD| header dependency tracking & Confirmed genuinely working — but specific to FreeBSD's own \verb|bsd.sys.mk|; absent from NetBSD's bundled set; MacPorts status unknown \\
\bottomrule
\end{longtable}
\end{center}

%===============================================================
\chapter{Tutorial: Building \texttt{System} and \texttt{Hello} From Scratch}
%===============================================================

This tutorial builds the same worked example used throughout the
design documents: two frameworks, \texttt{System} (no dependencies)
and \texttt{Hello} (depends on \texttt{System}), where \texttt{Hello}
contains an executable module \texttt{hello.m} and a shared-library
module \texttt{libgreet.m} that \texttt{hello.m} links
against\reqfoot{REQ-framework-contains-modules}.

\section{Step 1 --- Create the Workspace}

\begin{lstlisting}[language=bash]
$ mkdir -p myworkspace
$ cd myworkspace
$ cat > makefile <<'EOF'
.include <mk.workspace.mk>
PARENT_WS=
EOF
\end{lstlisting}

An empty \verb|PARENT_WS=| is valid — this workspace has no parents.

\section{Step 2 --- Create the \texttt{System} Framework}

\begin{lstlisting}[language=bash]
$ mkdir -p System/include System/local/include
$ cat > System/makefile <<'EOF'
.include <mk.framework.mk>
PREREQS=
EOF
\end{lstlisting}

The empty, but present, \verb|PREREQS=| is what makes \texttt{System}
auto-discoverable as a framework at all\reqfoot{REQ-prereqs-mandatory-even-empty-req}.

\section{Step 3 --- Create the \texttt{Hello} Framework}

\begin{lstlisting}[language=bash]
$ mkdir -p Hello/include Hello/local/include
$ cat > Hello/makefile <<'EOF'
.include <mk.framework.mk>
PREREQS=System
EOF
\end{lstlisting}

\section{Step 4 --- Create the \texttt{libgreet.m} Library Module}

\begin{lstlisting}[language=bash]
$ mkdir -p Hello/libgreet.m/src Hello/libgreet.m/include
$ cat > Hello/libgreet.m/makefile <<'EOF'
.include <mk.lib.mk>
LIB=greet
LIB_SHARED=YES
SHLIB_MAJOR=1
LIBS=
EOF
$ cat > Hello/libgreet.m/include/greet.h <<'EOF'
void greet(void);
EOF
$ cat > Hello/libgreet.m/src/greet.c <<'EOF'
#include "greet.h"
#include <stdio.h>
void greet(void) { printf("hello from libgreet\n"); }
EOF
\end{lstlisting}

Note \verb|greet.h| lives in the module's own \emph{private}
\verb|include/| — visible only to \verb|libgreet.m|'s own sources, not
even to sibling module \verb|hello.m|, regardless of the fact that
\verb|hello.m| will link against
it\reqfoot{REQ-three-level-header-visibility-req, REQ-libs-macro-link-only-headers-via-prereqs-req}.
If \texttt{hello.m} needed to \verb|#include "greet.h"| directly, it
would need to live in \texttt{Hello}'s \verb|local/include/| instead.

\section{Step 5 --- Create the \texttt{hello.m} Executable Module}

\begin{lstlisting}[language=bash]
$ mkdir -p Hello/hello.m/src
$ cat > Hello/hello.m/makefile <<'EOF'
.include <mk.prog.mk>
PROG=hello
LIBS=greet
EOF
$ cat > Hello/hello.m/src/main.c <<'EOF'
extern void greet(void);
int main(void) { greet(); return 0; }
EOF
\end{lstlisting}

\section{Step 6 --- Build}

\begin{lstlisting}[language=bash]
$ cd Hello
$ make
\end{lstlisting}

Behind the scenes: \verb|libgreet.m| is built first (it has no
dependency on \verb|hello.m|), producing \verb|libgreet.so.1| with a
\verb|libgreet.so| symlink; \verb|hello.m| is built second, since its
\verb|LIBS=greet| entry makes it depend on \verb|libgreet.m| having
been \emph{fully} built first\reqfoot{REQ-fused-build-per-module-full-dependency-order-req}.
Both modules' artifacts are then copied up into
\verb|Hello/build/<KEY>/{bin,lib}/|\reqfoot{REQ-aggregation-step-copies-outputs}.

\section{Step 7 --- Run It}

\begin{lstlisting}[language=bash]
$ cd ..    # back to the workspace root
$ make run PROGRAM=hello
hello from libgreet
\end{lstlisting}

\verb|make run| sets \verb|LD_LIBRARY_PATH| (or
\verb|DYLD_LIBRARY_PATH| on macOS) to include the workspace's
consolidated \verb|build/<KEY>/lib/| before executing, so
\verb|libgreet.so| is found without needing a real \verb|make install|
step first\reqfoot{REQ-run-target-req}.

\section{Step 8 --- Clean Up}

\begin{lstlisting}[language=bash]
$ make clean TARGET=all
\end{lstlisting}

This removes every generated artifact for every target present — the
whole \verb|build/| tree, any \verb|distrib/| packaging staging, and
all internal caches — leaving the workspace indistinguishable from a
fresh checkout\reqfoot{REQ-clean-removes-all-derived-artifacts-req}.

\vfill
\begin{center}
\rule{0.5\textwidth}{0.4pt}\\[0.3em]
\small
This document cites the discovery-driven design record maintained in
\texttt{project-model/}. Every footnoted alias (e.g.\
\texttt{REQ-framework-contains-modules}) corresponds to exactly one
file there, carrying its full rationale, decision history, and any
superseded predecessors.
\end{center}

\end{document}
